\section{Confiabilidade e expansão}

Diante disso, as operações podem ser mais seguras e eficientes, porém demandam menor intervenção do controle humano, enquanto a redundância e o uso de algoritmos avançados reduzem a probabilidade de erros, a evolução nessa área é constante.
A redução de controle remoto vindos de estações espaciais na Terra minimiza a latência nas respostas e aprimora a autonomia dos sistemas espaciais, ou seja, obtém-se maior confiabilidade em ambientes complexos e de microgravidade.
Além disso, a integração de sistemas de RVD/B com tecnologias autônomas e a fusão de sensores possibilitam missões cada vez mais complexas e ousadas.